% Problem-section file following ../_template.tex.
\section{Generalized Stein lemma for fully quantum channel resources}
\label{sec:problem-47}

\paragraph{Problem.}
Let $\mathcal N:\mathcal L(A)\to\mathcal L(B)$ be a finite-dimensional
quantum channel.  For each $n$, let $\mathfrak F_n$ be a nonempty compact,
convex, permutation-invariant set of channels from $A^{\otimes n}$ to
$B^{\otimes n}$, closed under tensor products and containing
$\mathcal R_\omega^{\otimes n}$ for one full-rank state $\omega$, where
$\mathcal R_\omega(X):=\operatorname{Tr}(X)\omega$.  Define
\begin{equation}
  D_{\rm ch}(\mathcal N^{\otimes n}\|\mathcal M_n)
  :=\sup_{\psi_{R_nA^n}}
  D\!\left(
    (\operatorname{id}_{R_n}\otimes\mathcal N^{\otimes n})(\psi)
    \middle\|
    (\operatorname{id}_{R_n}\otimes\mathcal M_n)(\psi)
  \right),
  \qquad R_n\simeq A^{\otimes n},
  \label{eq:p47-channel-relative-entropy}
\end{equation}
where the supremum in Eq.~\eqref{eq:p47-channel-relative-entropy} is over
density operators and $D$ is the quantum relative entropy.  The distance to
the free set is
\begin{equation}
  E_n^{\rm QQ}(\mathcal N\|\mathfrak F_n)
  :=\inf_{\mathcal M_n\in\mathfrak F_n}
  D_{\rm ch}(\mathcal N^{\otimes n}\|\mathcal M_n).
  \label{eq:p47-free-channel-distance}
\end{equation}
Equation~\eqref{eq:p47-free-channel-distance} is the $n$-use relative-entropy
distance to the free channel set.  For $\varepsilon\in(0,1)$, define the
optimal worst-case type-II error of a
parallel quantum-input/quantum-output test by
\begin{equation}
  \begin{aligned}
  \beta_{\varepsilon,n}^{\rm QQ}(\mathcal N\|\mathfrak F_n)
  :=\inf_{\substack{\psi_{R_nA^n},\ 0\leq Q\leq I\\
       \operatorname{Tr}[Q(\operatorname{id}_{R_n}\otimes
       \mathcal N^{\otimes n})(\psi)]\geq1-\varepsilon}}
  \ \sup_{\mathcal M_n\in\mathfrak F_n}
  \operatorname{Tr}\!\left[
       Q(\operatorname{id}_{R_n}\otimes\mathcal M_n)(\psi)
  \right].
  \end{aligned}
  \label{eq:p47-composite-type-two-error}
\end{equation}
Under what additional structural assumptions on $(\mathfrak F_n)_{n\geq1}$,
if any, do both limits exist and obey the fully quantum generalized Stein
identity
\begin{equation}
  \lim_{n\to\infty}-\frac1n\log_2
  \beta_{\varepsilon,n}^{\rm QQ}(\mathcal N\|\mathfrak F_n)
  =\lim_{n\to\infty}\frac1n
  E_n^{\rm QQ}(\mathcal N\|\mathfrak F_n)
  \qquad\text{for every }\varepsilon\in(0,1)?
  \label{eq:p47-fully-quantum-stein}
\end{equation}

\subsection*{Status}

Unsolved

\subsection*{Source}

The fully quantum formulation is implicit in the two generalized Stein
theorems for classical-input channels, both of which isolate their
classical-input structure from the unresolved quantum-input setting
\sourcecite{ref:p47-hayashi-yamasaki-cq}{HY25b},
\sourcecite{ref:p47-bergh-datta-khaitan}{BDK25}.

\subsection*{Progress}

\begin{itemize}
  \item For an i.i.d. resource state tested against admissible composite sets
  of free states, the generalized quantum Stein lemma identifies the
  fixed-error exponent with the regularized relative entropy of resource
  \sourcecite{ref:p47-hayashi-yamasaki-state}{HY25a}.

  \item Two independent works prove the channel analogue for
  classical--quantum channels.  In that setting the channel relative entropy
  reduces to
  \begin{equation}
    D_{\rm CQ}(\Phi\|\Psi)=\max_x D(\rho_x\|\sigma_x),
    \qquad \Phi:x\mapsto\rho_x,\quad\Psi:x\mapsto\sigma_x,
    \label{eq:p47-cq-divergence}
  \end{equation}
  and its pointwise structure supplies the additivity and minimax steps needed
  for a fixed-error theorem.  These steps apply to
  Eq.~\eqref{eq:p47-cq-divergence} but are unavailable in this form for QQ
  channels
  \sourcecite{ref:p47-hayashi-yamasaki-cq}{HY25b},
  \sourcecite{ref:p47-bergh-datta-khaitan}{BDK25}.

  \item The CQ proof explicitly states that analogous properties for fully
  quantum channels remain unclear.  Entangled quantum inputs introduce a
  reference system and a nontrivial order between input optimization and the
  worst-case free-channel optimization in
  Eq.~\eqref{eq:p47-composite-type-two-error}
  \sourcecite{ref:p47-hayashi-yamasaki-cq}{HY25b}.
\end{itemize}

\subsection*{References}

\begin{enumerate}[label={},leftmargin=4.8em,labelsep=0.5em]
  \footnotesize
  \item[\textup{[HY25a]}]\label{ref:p47-hayashi-yamasaki-state}
  M. Hayashi and H. Yamasaki,
  ``The Generalized Quantum Stein's Lemma and the Second Law of Quantum
  Resource Theories,'' \emph{Nature Physics} \textbf{21}, 1988--1993 (2025).
  \href{https://doi.org/10.1038/s41567-025-03047-9}%
       {doi:10.1038/s41567-025-03047-9};
  \href{https://arxiv.org/abs/2408.02722}{arXiv:2408.02722}.

  \item[\textup{[HY25b]}]\label{ref:p47-hayashi-yamasaki-cq}
  M. Hayashi and H. Yamasaki,
  ``Generalized Quantum Stein's Lemma for Classical-Quantum Dynamical
  Resources,'' arXiv preprint (2025).
  \href{https://arxiv.org/abs/2509.07271}{arXiv:2509.07271}.

  \item[\textup{[BDK25]}]\label{ref:p47-bergh-datta-khaitan}
  B. Bergh, N. Datta, and A. Khaitan,
  ``Generalized Quantum Stein's Lemma and Reversibility of Quantum Resource
  Theories for Classical-Quantum Channels,'' arXiv preprint (2025).
  \href{https://arxiv.org/abs/2509.13280}{arXiv:2509.13280}.
\end{enumerate}

\subsection*{Comment}

The state theorem and the CQ-channel theorems do not imply
Eq.~\eqref{eq:p47-fully-quantum-stein} for genuinely quantum inputs.  An
adaptive quantum-comb version would be a further problem and is not included
in the present statement.

\subsection*{Tag}

Quantum hypothesis testing; Quantum resource theories; Quantum channels;
Superchannels

\subsection*{ID}

\texttt{op\_2579e084f37ac18c}
